%% Complete editorial counting proof; original-source acquisition is separate.
\section{The pigeonhole counting argument}

\begin{annexitem}{annex:p7.complexity.1}{thm:p07-haken-resolution-lower-bound}{An explicit exponential count for pigeonhole resolution}
\begin{proof}
We prove the more explicit bound
\[
 L\ \geq\ (16/15)^{\lfloor n/4\rfloor}\qquad(n\geq4),
\]
where \(L\) is the number of clause occurrences in the given DAG.
Rows below are pigeons and columns are holes, as in
Definition~\ref{def:p07-pigeonhole-cnf}. Clauses retain the literal-set
convention of Sheet~\ref{sheet:resolution-cpl}. In particular, no
functionality clauses or weakening rule are added.

\paragraph{Critical assignments and a one-step bound.}
A critical assignment matches \(n\) of the \(n+1\) pigeons bijectively
to the \(n\) holes, and leaves the remaining pigeon unmatched. It makes
precisely those matching variables true. Every collision input is true;
the unique false input is the unmatched pigeon's row clause.
At a false resolvent, follow the parent whose pivot literal is false.
Every other literal of that parent is also a literal of the resolvent,
so the selected parent is false. This is the routing proved in
Theorem~\ref{thm:p07-haken-critical-count-routing}.

Order all occurrences topologically, with both parents earlier than their
resolvent. Routing to parents strictly decreases the index and ends at
the unique false input. Write
\[
 p_i(C)=\#\{j:x_{ij}\in C\}
\]
for the number of positive literals in pigeon row \(i\) of clause \(C\).
On a step from a resolvent to its selected parent this number increases
by at most one: only the removed pivot literal can be in the parent
without also being in the resolvent. It may decrease by more than one;
no monotonicity in the other direction is assumed. If the unmatched
pigeon of a falsifying critical assignment is \(i\), the false clause
has no negative literal in row \(i\), because every variable there is
false.

\paragraph{Assigning partial matchings to occurrences.}
Put \(k=\lfloor n/4\rfloor\), so \(1\leq k\) and \(2k<n\).
Let \(\mathcal M_{n,k}\) be the family of all injections \(f\) from
a \(k\)-element subset of the \(n+1\) pigeons into the \(n\) holes.
A critical assignment extends \(f\) when it makes all \(x_{i,f(i)}\),
for \(i\) in the domain of \(f\), true. Every \(f\) has such an extension:
choose an unmatched pigeon outside its domain and biject the remaining
unmatched pigeons with the unused holes.

For any such extension with unmatched pigeon \(u\), its route from the
empty clause to the row input has positive-row counts starting at zero
and ending at \(n\). The one-step bound forces a count of exactly \(2k\)
at some occurrence. Call an occurrence \emph{eligible for \(f\)} when
some critical assignment extending \(f\) falsifies it and has unmatched
pigeon \(u\) with \(p_u(C)=2k\). Choose \(C_f\) to be the earliest
eligible occurrence in the fixed topological order. This choice is
well-defined on the original DAG, without unfolding shared nodes.

The minimality of \(C_f\) has the following useful consequence.
If a critical assignment extending \(f\) falsifies \(C_f\) and leaves
\(q\) unmatched, then
\[
 p_q(C_f)\geq2k.                                                \tag{*}
\]
Indeed, if its count at \(C_f\) were smaller than \(2k\), route this
assignment from \(C_f\) to its false row input. The count ends at \(n\)
and increases by at most one per step. An earlier occurrence would
therefore have count exactly \(2k\) and would be eligible for \(f\),
contrary to the choice of \(C_f\).

\paragraph{At least \(k+1\) constrained rows.}
Fix \(f\) and a witness \(v\) to the eligibility of \(C_f\), with unmatched
pigeon \(u\). There are \(2k\) holes whose positive literals occur in
row \(u\) of \(C_f\), and \(k\) holes in the image of \(f\).
At least \(n-3k\geq k\) holes belong to neither set.
For each such hole \(j\), let \(q\) be the pigeon matched to \(j\) by
\(v\). The pigeons obtained from distinct holes are distinct,
are different from \(u\), and do not belong to the domain of \(f\).

Move the matching edge at hole \(j\) from \(q\) to \(u\). The resulting
assignment \(v'\) is critical, still extends \(f\), and leaves \(q\)
unmatched. Only \(x_{u,j}\) and \(x_{q,j}\) change:
the first from false to true and the second from true to false.
The positive literal \(x_{u,j}\) is absent from \(C_f\) by the choice
of \(j\). If \(\neg x_{q,j}\) occurs, row \(q\) contains a negative
literal. If it does not occur, neither changed variable makes a
literal of \(C_f\) true, so \(v'\) still falsifies \(C_f\).
Consequence \((*)\) then gives \(p_q(C_f)\geq2k\).
Together with row \(u\), this proves that \(C_f\) has at least \(k+1\)
distinct rows each of which contains a negative literal or at least
\(2k\) positive literals.

For any occurrence \(C\) in the image of \(f\mapsto C_f\), fix a set
\(A\) of exactly \(k+1\) such rows. The set is chosen from \(C\) once,
independently of which \(f\)'s are subsequently counted at that occurrence.
Suppose a critical assignment falsifies \(C\). If row \(i\in A\)
contains \(\neg x_{i,j}\), that assignment must match \(i\) to \(j\).
Thus a partial matching contained in the assignment has at most one
possible image for \(i\) when it includes that pigeon. If the row instead
has at least \(2k\) positive literals, all their holes are forbidden
as its image. In either case there are at most \(n-2k\) possible images.
The assertion remains valid when several negative literals make
falsification impossible; such a clause simply has no matching to count.
Here \(1\leq n-2k\), as required for the first case.

\paragraph{An explicit bound on each fiber.}
The total number of partial matchings is
\[
 H=\#\mathcal M_{n,k}=\binom{n+1}{k}\frac{n!}{(n-k)!}.
\]
For a fixed image occurrence \(C\), let \(i\) of a partial matching's
\(k\) pigeons lie in its fixed set \(A\). The number of choices of
domain with this property is
\[
 d_i=\binom{k+1}{i}\binom{n-k}{k-i},\qquad 0\leq i\leq k.
\]
For a fixed such domain, the \(i\) constrained pigeons have at most
\((n-2k)^i\) image choices if the requirement that their images be
distinct is ignored. For each valid injective choice among them, the
other \(k-i\) labeled pigeons have at most
\((n-i)!/(n-k)!\) injective image choices. Multiplying these upper
bounds counts every valid matching, possibly with surplus choices;
it does not treat different images as independent random events.
Consequently every fiber has size at most
\[
 G=\sum_{i=0}^{k}d_i(n-2k)^i\frac{(n-i)!}{(n-k)!}.
\]
The fibers partition \(\mathcal M_{n,k}\), and there are at most \(L\)
image occurrences, hence \(H\leq LG\).
By counting the \(k\)-element domains according to their intersection
with \(A\), \(\sum_{i=0}^{k}d_i=\binom{n+1}{k}\). It follows that
\[
 \frac GH
 =\frac{\displaystyle\sum_{i=0}^{k}d_i
          \frac{(n-2k)^i(n-i)!}{n!}}
        {\displaystyle\binom{n+1}{k}}.
\]
Each factorial ratio is at most
\[
 \left(\frac{n-2k}{n-k+1}\right)^i\leq(3/4)^i.
\]
For the last inequality write \(n=4k+r\), \(0\leq r\leq3\):
\(4(2k+r)\leq3(3k+r+1)\) follows from \(r\leq k+3\).
The empty product when \(i=0\) is one.

\paragraph{The exponential estimate, including its averaging lemma.}
We use the following elementary inequality. Let \(w_1,\ldots,w_N\) be
nonnegative. Write \(e_\ell(w)\) for the sum of the products over
all \(\ell\)-element subsets of the indices, with \(e_0=1\).
Then, for \(1\leq\ell\leq N\),
\[
 \frac{e_\ell(w)}{\binom N\ell}
 \leq\left(\frac{w_1+\cdots+w_N}{N}\right)^\ell.             \tag{**}
\]
For \(\ell=1\) this is equality. For \(\ell\geq2\), keep all but
two coordinates \(a,b\) fixed. Separating subsets by how many of these
coordinates they contain gives
\[
 e_\ell(w)=ab\,e_{\ell-2}(z)+(a+b)e_{\ell-1}(z)+e_\ell(z),
\]
where \(z\) lists the other coordinates and an elementary sum beyond
its range is zero. Replacing \(a,b\) by their common average preserves
their sum and increases their product, so it cannot decrease
\(e_\ell\). Repeatedly average a largest and a smallest coordinate.
Their sum remains fixed. If their common overall mean is \(\mu\),
the quantity \(\sum_j(w_j-\mu)^2\) decreases by
\((a-b)^2/2\) at each step. If it had a positive limiting value,
the largest-minus-smallest difference would stay bounded below,
since \(\sum_j(w_j-\mu)^2\leq N(\max w_j-\min w_j)^2\).
That would force a fixed positive decrease indefinitely, a contradiction.
Thus the coordinates converge to \(\mu\). Continuity of \(e_\ell\)
now bounds its initial value by \(\binom N\ell\mu^\ell\), proving
\((**)\).

Apply \((**)\) with \(N=n+1\), \(\ell=k\), weight \(3/4\) on the
\(k+1\) indices in \(A\), and weight one on all other indices.
Grouping the products by their intersection with \(A\) gives
\[
 \begin{aligned}
 \frac GH
 &\leq\frac{\sum_{i=0}^{k}d_i(3/4)^i}{\binom{n+1}{k}}\\
 &\leq\left(1-\frac{k+1}{4(n+1)}\right)^k
 \leq(15/16)^k.
 \end{aligned}
\]
The last inequality uses \(4(k+1)\geq n+1\).
Since \(H\leq LG\), we obtain \(L\geq(16/15)^k\).
For \(n\geq8\), \(k\geq n/8\), so
\(c=(16/15)^{1/8}>1\) is one constant satisfying the theorem.

Every object counted is a partial matching or an occurrence in the
original annotated DAG. Shared clauses are never unfolded into a tree;
equal literal sets at different occurrences remain distinct occurrences.
Tautological clauses have no falsifying critical assignment and cannot
be selected. Thus the proof covers exactly the registered resolution
rule, including its permission for tautological intermediate clauses,
and the nonfunctional input family of the stated theorem.
\end{proof}
\end{annexitem}
